\documentclass[10pt]{article}
\usepackage[margin=0.82in,top=0.65in,bottom=0.55in]{geometry}
\usepackage{amsmath,amssymb,amsthm}
\usepackage[protrusion=true,expansion=false]{microtype}
\usepackage{xcolor}
\usepackage[most]{tcolorbox}
\usepackage{titlesec}
\usepackage{enumitem}
\definecolor{acc}{HTML}{1B6B58}
\definecolor{grey}{HTML}{5C6763}
\titleformat{\section}{\normalfont\bfseries}{}{0em}{}
\titlespacing*{\section}{0pt}{1.1em}{0.3em}
\titleformat{\subsection}{\normalfont\bfseries\small}{}{0em}{}
\titlespacing*{\subsection}{0pt}{0.8em}{0.2em}
\newtcolorbox{idea}{enhanced, sharp corners, colback=white, colframe=acc!35,
  boxrule=0.5pt, left=0.7em, right=0.6em, top=0.3em, bottom=0.3em,
  before skip=0.4em, after skip=0.45em}
\newcommand{\Q}[2]{\medskip\noindent{\color{acc}\textbf{#1}}\ {\scriptsize\color{grey}\textsc{#2}}\quad}
\newcommand{\R}{\mathbb{R}}
\newcommand{\E}{\mathbb{E}}
\newcommand{\Prb}{\mathbb{P}}
\newcommand{\Var}{\operatorname{Var}}
\newcommand{\Cov}{\operatorname{Cov}}
\newcommand{\Tr}{\operatorname{tr}}
\newcommand{\Nor}{\mathcal{N}}
\newcommand{\T}{^{\!\top}}
\newcommand{\ip}[2]{\langle #1,#2\rangle}
\newcommand{\KL}[2]{D(#1\,\|\,#2)}
\newcommand{\Ber}{\mathrm{Bernoulli}}
\newcommand{\Poi}{\mathrm{Poisson}}
\newcommand{\TV}{d_{\mathrm{TV}}}
\pagestyle{empty}
\setlength{\parskip}{0.26em}
\setlength{\parindent}{0pt}

\begin{document}

{\footnotesize\color{acc}\textsc{ORF 526 \quad$\cdot$\quad Solutions to all exercises, Lectures 1--6 \quad$\cdot$\quad instructor copy}}

\vspace{0.3em}
{\large\bfseries Solutions}
\vspace{0.2em}

\hrule height 0.5pt
\vspace{0.5em}

Numbering follows the lecture notes. For \textsc{type c} exercises the idea is boxed; the rest are routine once the idea is in hand, which is what the types mean.

\section{Lectures 1 and 2}
\subsection{\S1.8$^\ast$ Entropy, and why the Gaussian maximizes it}

\Q{1.13}{A} $H(q)=-q\log q-(1-q)\log(1-q)$: symmetric about $q=\tfrac12$, zero at the endpoints, maximal at $\tfrac12$ with value $\log2$. Uniform on $n$ points: $\log n$.

\Q{1.14}{B} By Jensen for the concave $\log$,
$H(p)=\E[\log\tfrac1{p(X)}]\le\log\E[\tfrac1{p(X)}]=\log\#\{x:p(x)>0\}\le\log n$.
Equality in Jensen forces $1/p(X)$ a.s.\ constant, so $p$ is uniform on its support; equality in the second step forces full support.

\Q{1.15}{A} $\log a$; $1-\log\lambda$; $\tfrac12\log(2\pi e\sigma^2)$. Uniform on $[0,a]$ with $a<1$ has $h=\log a<0$. Discrete entropy cannot be negative because $-p\log p\ge0$ for $p\in(0,1]$; a \emph{density} may exceed $1$, which is exactly what breaks that.

\Q{1.16}{A} $p_Y(y)=|a|^{-1}p_X(y/a)$; substituting $x=y/a$,
$h(Y)=-\int p_X(x)[\log p_X(x)-\log|a|]dx=h(X)+\log|a|$.

\Q{1.17}{B} $-\KL{p}{q}=\int p\log\tfrac qp\le\log\int q=0$ by Jensen, with equality iff $q/p$ is a.s.\ constant, hence $\equiv1$. Under $X\mapsto aX$ both densities gain the same factor $|a|^{-1}$ and the same substitution is made, so $\KL{p}{q}$ is unchanged.

\Q{1.18}{B} $\log q(x)=-\tfrac n2\log(2\pi)-\tfrac12\log\det\Sigma-\tfrac12(x-\mu)\T\Sigma^{-1}(x-\mu)$ is quadratic in $x$, so $\int p\log q$ is a fixed linear functional of $\int p$, $\int x_ip$, $\int x_ix_jp$. Matching mean and covariance therefore gives $\int p\log q=\int q\log q=-h(q)$, and $0\le\KL{p}{q}=-h(p)+h(q)$. For $h(q)$: $\E[(X-\mu)\T\Sigma^{-1}(X-\mu)]=\Tr(\Sigma^{-1}\Sigma)=n$, so $h(q)=\tfrac12\log((2\pi e)^n\det\Sigma)$.

\Q{1.19}{C}
\begin{idea}
$(x,y)\mapsto\big(\tfrac{x+y}{\sqrt2},\tfrac{x-y}{\sqrt2}\big)$ is \textbf{orthogonal}: Jacobian $1$, so it preserves joint differential entropy.
\end{idea}
With $U=\tfrac{X+Y}{\sqrt2}$, $V=\tfrac{X-Y}{\sqrt2}$: $h(U,V)=h(X,Y)=2h(X)$ by independence. Subadditivity $h(U,V)\le h(U)+h(V)$ is Exercise 1.17 applied to the joint density against the product of its marginals. Symmetry gives $(X,Y)\overset{d}{=}(X,-Y)$, so $V\overset{d}{=}U$. Chaining, $2h(X)\le2h(U)$.

\Q{1.20}{C}
\begin{idea}
The Gaussian proof never used Gaussianity. It used that $\log q$ is an \textbf{affine combination of the constrained functions}, so $\int p\log q$ is fixed by the constraints alone.
\end{idea}
$\log q=\sum_m\lambda_mT_m-\log Z$, so for any feasible $p$, $0\le\KL{p}{q}=-h(p)-\sum_m\lambda_mc_m+\log Z$ --- and the last two terms do not depend on $p$. Taking $p=q$ evaluates them as $h(q)$; subtracting gives $h(p)\le h(q)$, equality iff $p=q$. (No integrability of $h(p)$ need be assumed: $-\infty$ makes it trivial, $+\infty$ is excluded by $\KL{p}{q}\ge0$.)

Instances: no constraint on $[0,a]$ gives $q$ constant; $T_1(x)=x$ on $[0,\infty)$ gives the exponential; $T$'s $=x_i,x_ix_j$ give $\Nor(\mu,\Sigma)$.

\emph{Why the assumption is not free.} On $S=\R$ with only the mean constrained, $e^{\lambda x}$ is never integrable, and correctly so: $\Nor(0,\sigma^2)$ is feasible for every $\sigma$ with $h=\tfrac12\log(2\pi e\sigma^2)\to\infty$. \textbf{``$q$ is a density'' is exactly what rules this out}, and a student who assumes existence has assumed the content.

\subsection{\S2.5 Bayesian linear models, and \S2.6$^\ast$}

\Q{2.9}{A} $p(y\mid w)=(2\pi\sigma^2)^{-n/2}e^{-\|y-Hw\|^2/2\sigma^2}$, so $-\log p=\frac{1}{2\sigma^2}\|y-Hw\|^2+\frac n2\log(2\pi\sigma^2)$ and the minimizer over $w$ is the least-squares fit. Choosing squared error is asserting that the noise is Gaussian, homoscedastic and independent across observations; most people choosing it are not asserting anything, which is the point.

\Q{2.10}{A} The posterior density is $\propto\exp\big(-\frac{1}{2\sigma^2}\|y-Hw\|^2-\frac{1}{2\tau^2}\|w\|^2\big)$, so its maximizer minimizes $\|y-Hw\|^2+\lambda\|w\|^2$ with $\lambda=\sigma^2/\tau^2$. Large $\lambda$ means noisy data against a confident prior, so heavy shrinkage; $\tau\to0$ gives $\widehat w=0$ and $\tau\to\infty$ gives least squares.

\Q{2.11}{B} $(w,y)$ is a linear image of $(w,\xi)$, hence jointly Gaussian, with
$\Sigma_{ww}=\tau^2I_p$, $\Sigma_{wy}=\E[w(Hw+\sigma\xi)\T]=\tau^2H\T$, $\Sigma_{yy}=\tau^2HH\T+\sigma^2I_n$.
Corollary 2.5 then gives the two displays directly. $\Cov(w\mid y)$ contains no $y$: the data changes your \emph{estimate}, never your \emph{uncertainty}. Operationally, a surprising $y$ moves the posterior mean and leaves the error bars exactly where they were --- special to the Gaussian model, and false in general.

\Q{2.12}{B} $\Cov(X,Y)=\E[(Y^2-1)Y]=\E Y^3=0$, so the best linear predictor is $0$ with error $\Var X=\Var(Y^2)=2$. The predictor $f(Y)=Y^2-1$ has error $0$. What fails is \textbf{joint} Gaussianity of $(X,Y)$ --- $X$ is not Gaussian at all --- so Theorem 2.3 does not apply, and the conditional mean is not linear.

\Q{2.13}{C}
\begin{idea}
$H\T(HH\T+\lambda I_n)=(H\T H+\lambda I_p)H\T$ --- one line of expansion. Multiply by the two inverses.
\end{idea}
Both sides equal $H\T HH\T+\lambda H\T$. Left-multiplying by $(H\T H+\lambda I_p)^{-1}$ and right-multiplying by $(HH\T+\lambda I_n)^{-1}$ gives the identity. The right-hand form inverts an $n\times n$ matrix, so the cost is $O(n^3)$; this is the kernel trick, and why Gaussian process regression is feasible.

As $\tau\to\infty$ we have $\lambda\to0$, and in the form $H\T(HH\T+\lambda I_n)^{-1}y$ the limit is benign because $HH\T$ is $n\times n$ and invertible: $\E[w\mid y]\to H\T(HH\T)^{-1}y$, the minimum-norm solution of $Hw=y$. In the other form the limit is $0\cdot\infty$ --- which is what makes this a \textsc{c} and not a computation.

\section{Lectures 3 and 4}
\subsection{\S1.7$^\ast$ The proof we did not give}

\Q{1.9}{A} Differentiating $\varphi_X(t)=\E[e^{itX}]$ under the integral $k$ times brings down $(iX)^k$. For $\varphi=e^{-t^2/2}$: $\varphi'=-t\varphi$, $\varphi''=(t^2-1)\varphi$, so $\varphi''(0)=-1=i^2\E g^2$. The series $e^{-t^2/2}=1-\tfrac{t^2}{2}+\tfrac{t^4}{8}-\cdots$ has $t^4$ coefficient $\tfrac18$, so $\varphi^{(4)}(0)=4!/8=3=i^4\E g^4$.

\Q{1.10}{B} With $\theta=tX$, $\E[e^{itX}]=1-\tfrac{t^2}{2}+r(t)$ where $|r(t)|\le\E\min(|t|^3|X|^3/6,\ t^2X^2)$. Divide by $t^2$: the integrand is bounded by $X^2$, which is integrable, and tends to $0$ pointwise, so dominated convergence gives $r(t)/t^2\to0$. \emph{The second branch of the minimum is what supplies the dominating function without a third moment.}

\Q{1.11}{B} With $a=1+z_n/n$ and $b=e^{z_n/n}$, both are $1+O(1/n)$ so $|a|,|b|\le M$ for some fixed $M$, and $|a-b|=O(1/n^2)$; hence $|a^n-b^n|\le nM^{n-1}|a-b|\to0$ while $b^n=e^{z_n}\to e^z$. Applying this with $z_n=n(\varphi_{X_1}(t/\sqrt n)-1)\to-t^2/2$ by Exercise 1.10 gives $\varphi_{S_n}(t)\to e^{-t^2/2}$, and L\'evy finishes. \textbf{Independence} is what gives $\varphi_{S_n}=\prod_i\varphi_{X_i}(t/\sqrt n)$; \textbf{identical distribution} only lets us write one factor $n$ times.

\Q{1.12}{A} $\varphi_B(t)=(1-p)+pe^{it}=1+p(e^{it}-1)$. So $\varphi_{N_n}=[1+\tfrac\lambda n(e^{it}-1)]^n\to\exp(\lambda(e^{it}-1))$ by Exercise 1.11, and $\sum_ke^{-\lambda}\frac{\lambda^k}{k!}e^{itk}=e^{-\lambda}e^{\lambda e^{it}}$ is the same thing.

\Q{1.13}{B} $X_{n,i}$ takes $1-\lambda/n$ with probability $\lambda/n$ and $-\lambda/n$ otherwise, so $s_n^2=n\cdot\frac\lambda n(1-\frac\lambda n)=\lambda(1-\frac\lambda n)$. For $\varepsilon<1/\sqrt\lambda$ the threshold $\varepsilon s_n\to\varepsilon\sqrt\lambda<1$ eventually falls below the large value, so $L_n(\varepsilon)\to n\cdot\frac\lambda n\cdot1/\lambda=1$.

At $\lambda=3$, $\varepsilon=\tfrac12$: $n=10$ gives $s_n=1.449$ and $\varepsilon s_n=0.7246$, while the large value is $1-\tfrac3{10}=0.7<0.7246$ --- nothing is caught, so $L_{10}=0$ \textbf{exactly}. Then $0.970$ and $1.000$.

\Q{1.14}{C}
\begin{idea}
A Gaussian characteristic function $e^{-\sigma^2t^2/2}$ \textbf{never vanishes}. So exhibiting one zero of the limit is a complete proof that the limit is not Gaussian.
\end{idea}
$s_n^2=n\cdot1+n=2n$. The large summand has $|X_{n,n+1}|=\sqrt n>\varepsilon\sqrt{2n}=\varepsilon s_n$ exactly when $\varepsilon<1/\sqrt2$, and it alone contributes $n/2n=\tfrac12$ to $L_n$. The first $n$ summands give $s_n^{-1}\sum_{i\le n}X_{n,i}\Rightarrow\Nor(0,\tfrac12)$ with transform $e^{-t^2/4}$; the last gives $\pm1/\sqrt2$ with transform $\cos(t/\sqrt2)$; they are independent, so the product is the limit. It vanishes at $t=\pi/\sqrt2$.

\Q{1.15}{C}
\begin{idea}
Write $\E f(X_\Sigma)=\frac{1}{(2\pi)^n}\int\hat f(t)\,e^{-\frac12t\T\Sigma t}\,dt$, so that $\Sigma$ appears \textbf{only in the exponent}.
\end{idea}
That identity is Fourier inversion plus Definition 1.2. Differentiating in $s$ under the integral brings down $-\tfrac12\sum_{ij}\dot\Sigma_{ij}t_it_j$. Since $\widehat{\partial_i\partial_jf}(t)=(it_i)(it_j)\hat f(t)=-t_it_j\hat f(t)$, this is $\tfrac12\sum_{ij}\dot\Sigma_{ij}\widehat{\partial_i\partial_jf}(t)$, and inverting the transform term by term gives the claim. Integrating over $\Sigma(s)=(1-s)\Sigma^0+s\Sigma^1$, for which $\dot\Sigma=\Sigma^1-\Sigma^0$, gives the second display.

Through the density instead one is differentiating $\Sigma^{-1}$ and $\det\Sigma$: the same answer, a page later.

\emph{Check.} In $\R^2$ with unit variances and correlation $\rho$, $\frac{d}{d\rho}\E[X_1^2X_2^2]=\E[4X_1X_2]=4\rho$, against Wick's $\frac{d}{d\rho}(1+2\rho^2)=4\rho$.

\Q{1.16}{C}
\begin{idea}
$S_n$ has an \textbf{atom}. A distribution function with a jump of size $a$ is at least $a/2$ from any continuous one.
\end{idea}
If $F(z_0)-F(z_0^-)=a$ and $\Phi$ is continuous then $|F(z_0)-\Phi(z_0)|+|F(z_0^-)-\Phi(z_0^-)|\ge a$, so the sup is $\ge a/2$. For $X_i=\pm1$ and $n$ even, $\Prb(S_n=0)=\binom{n}{n/2}2^{-n}\sim\sqrt{2/\pi n}$, giving $\sup_z|\Prb(S_n\le z)-\Phi(z)|\ge c\,n^{-1/2}$. Lattice summands were slowest in \S1.5 for exactly this reason: their atoms cannot be smoothed, and they are of the same order as the whole error budget.

\subsection{\S2.5$^\ast$ Where independence was spent}

\Q{2.10}{A} $\E F(W,W')=\E F(W',W)$ by exchangeability, and $=-\E F(W,W')$ by antisymmetry; so it is zero.

\Q{2.11}{B} $F(x,y)=(y-x)(f(x)+f(y))$ is antisymmetric, so $\E[(W'-W)(f(W)+f(W'))]=0$. Writing $f(W')=f(W)+(W'-W)f'(W)+R$ with $|R|\le\tfrac{\|f''\|_\infty}{2}(W'-W)^2$,
\[
0=2\E[(W'-W)f(W)]+\E[(W'-W)^2f'(W)]+\E[(W'-W)R].
\]
The first term is $-2a\E[Wf(W)]$ by the Stein-pair property and conditioning on $W$. Rearranging gives the display, with the remainder bounded by $\tfrac{\|f''\|_\infty}{4a}\E|W'-W|^3$.

\Q{2.12}{B} Swapping $X_I\leftrightarrow X_I'$ is measure preserving, so $(W,W')$ is exchangeable. Since $\E[X_I'\mid X_1,\dots,X_n]=0$ and $\E[X_I\mid X_1,\dots,X_n]=\frac1n\sum_iX_i$,
$\E[W'-W\mid X_1,\dots,X_n]=-\frac{1}{\sqrt n}\cdot\frac1n\sum_iX_i=-W/n$, and the tower property gives $a=1/n$.
Also $(W'-W)^2=(X_I-X_I')^2/n$, whose conditional expectation is $\frac1{n^2}\sum_i(X_i^2+1)$. Dividing by $2a=2/n$ gives $\tfrac12(1+\frac1n\sum_iX_i^2)$, whose distance from $1$ is $\tfrac12|1-\frac1n\sum X_i^2|$, at most $\tfrac12\sqrt{\Var(X_1^2)/n}$ in $L^1$. The remainder is $\tfrac{n\|f''\|_\infty}{4}n^{-3/2}\E|X_1-X_1'|^3$. Both are $O(n^{-1/2})$.

\emph{Where independence entered:} in the two conditional expectations, and nowhere else. The product formula $\E[X_if(W^{(i)})]=\E[X_i]\E[f(W^{(i)})]$ never appeared.

\Q{2.13}{A} $\E W=n\cdot\frac1n=1$; $\E[I_iI_j]=\frac{1}{n(n-1)}$ for $i\ne j$ gives $\E W^2=1+1=2$, so $\Var W=1$. And $\Prb(I_1=I_2=1)=\frac{1}{n(n-1)}\ne\frac1{n^2}$. What fails is \textbf{independence of the $I_i$}, which Theorem 2.7 requires; the theorem therefore says nothing, and a true conclusion reached from unsatisfied hypotheses is still not proved.

\Q{2.14}{C}
\begin{idea}
Force $i$ to be a fixed point by \textbf{composing with the transposition that sends $i$ to itself}: $\sigma=\pi\circ(i\;\;\pi^{-1}(i))$.
\end{idea}
Then $\sigma(i)=\pi(\pi^{-1}(i))=i$, and the map is a bijection from $\{\pi:\pi(i)\ne i\}$ onto itself composed with a relabelling, so $\sigma$ is uniform among permutations fixing $i$. Writing $j=\pi^{-1}(i)$: $i$ becomes fixed (gain $1$), and $j$ becomes fixed exactly when $\pi(i)=j$, i.e.\ when $\pi$ has a $2$-cycle at $i$ (gain $2$). Hence $V_i=W(\sigma)-1\in\{W,W+1\}$, and
\[
\E|W-V_i|\le\Prb(\pi(i)=i)+\Prb(\text{$2$-cycle at }i)=\tfrac1n+\tfrac1n .
\]
Chen--Stein with $p_i=1/n$, $\lambda=1$ gives $\TV(W,\Poi(1))\le n\cdot\frac1n\cdot\frac2n=\frac2n$.

At $n=10$ the bound is $0.2$ and the truth is $2.3\times10^{-8}$: wrong by about seven orders of magnitude. What it has that inclusion--exclusion does not is \textbf{robustness} --- it never used that $\pi$ was a permutation beyond the coupling, and it is the only one of the three routes in this course that produces a bound at all.

\section{Lectures 5 and 6}
\subsection{\S1 Wick's theorem}

\Q{1.6}{A} The first point has $2q-1$ choices of partner; deleting that block leaves $2q-2$ points, so the count is $(2q-1)!!$, giving $1,3,15,105$. Wick with $n=1$, $\Sigma=1$ makes every product of covariances equal $1$, so $\E[g^{2q}]$ is that count; $\E[g^{20}]=19!!$.

\Q{1.7}{B} $\E[X\T AX]=\sum_{ij}A_{ij}\Sigma_{ij}=\Tr(A\Sigma)$. Wick gives $\E[X_iX_jX_kX_l]=\Sigma_{ij}\Sigma_{kl}+\Sigma_{ik}\Sigma_{jl}+\Sigma_{il}\Sigma_{jk}$, so $\E[(X\T AX)^2]=(\Tr A\Sigma)^2+2\Tr(A\Sigma A\Sigma)$, the last two pair partitions being equal by symmetry of $A$ and $\Sigma$. Subtracting the squared mean, $\Var(X\T AX)=2\Tr((A\Sigma)^2)$; with $A=I$, $\Var\|X\|^2=2\Tr\Sigma^2$.

\subsection{\S2.9$^\ast$ Counting the walks}

\Q{2.12}{A} $\frac1N\E\Tr M^2=\frac1N\sum_{i,j}\E M_{ij}^2=\frac1N[(N^2-N)\tfrac1N+N\tfrac2N]=1+\frac1N$. Odd moments vanish with no computation: $\Tr M^{2k+1}$ is a sum of products of an odd number of centred Gaussians, and Wick pairs factors two at a time.

\Q{2.13}{B} The $r$th factor of a term is the matrix entry of the $r$th step of the closed walk $i_1\to\cdots\to i_{2k}\to i_1$. Since $\E[M_{ab}M_{cd}]=0$ unless $\{a,b\}=\{c,d\}$, a block joining two steps on different edges kills the term. So the $2k$ steps are matched into $k$ pairs, each pair on one edge: at most $k$ distinct edges, each used at least twice, hence exactly twice in the leading terms.

\Q{2.14}{B} A closed walk has a connected edge graph, and a connected graph on $v$ vertices has $e\ge v-1$; with $e\le k$ this gives $v\le k+1$, equality forcing $e=k$ and $e=v-1$, i.e.\ a tree with $k$ edges each traversed twice. Counting: $\sim N^v$ index tuples, $k$ covariances of size $1/N$, and a $1/N$ in front, so the contribution is $N^{v-k-1}$, which vanishes unless $v=k+1$.

\Q{2.15}{C}
\begin{idea}
Condition on the block containing the point $1$. Non-crossing forces the interval it spans to be self-contained, splitting the problem into two \textbf{independent} problems of the same kind.
\end{idea}
If that block is $\{1,b\}$, a block joining a point inside $(1,b)$ to one outside would cross it, so the $b-2$ interior points pair among themselves and $b=2j+2$ is even; the rest splits into $2j$ inside and $2(n-1-j)$ outside. Hence $C_n=\sum_{j<n}C_jC_{n-1-j}$, $C_0=1$: $1,1,2,5,14,42$. Multiplying by $w^n$ and summing turns the convolution into a product, $C-1=wC^2$; the root finite at $w=0$ is $C(w)=\frac{1-\sqrt{1-4w}}{2w}$, whose coefficients are $\frac{1}{n+1}\binom{2n}{n}$.

\Q{2.16}{C}
\begin{idea}
Forget where the walk is. Record only whether each step goes somewhere \textbf{new} ($+1$) or comes back ($-1$).
\end{idea}
On a tree each edge is traversed twice, and a first traversal necessarily lands on a vertex not yet visited, since a tree has no cycles. So $+1$ marks first traversals, the partial sum after $r$ steps is the number of edges opened but not closed --- the current depth, hence $\ge0$ --- and the walk is closed, so the total is $0$. That is a Dyck path. Reading a Dyck path left to right and opening a new edge at each $+1$ reconstructs the tree and the traversal, so the map is a bijection, and there are $C_k$ shapes by Exercise 2.15.

A contributing walk is a shape plus an injective labelling of its $k+1$ vertices, so the count is exactly $C_k\,N(N-1)\cdots(N-k)$. Multiplying by $N^{-k}$ and $N^{-1}$ gives $C_k(1+O(1/N))$.

\Q{2.17}{C}
\begin{idea}
The two $\delta$'s in the real covariance come from $M_{ij}=M_{ji}$ --- there are \textbf{two ways} to glue a pair of entries. Hermitian symmetry leaves one.
\end{idea}
With $\E[H_{ab}H_{cd}]=\frac1N\delta_{ad}\delta_{bc}$, the first block of the crossing pair partition forces $i_1=i_4$ and $i_2=i_3$, and the second forces $i_2=i_1$ and $i_3=i_4$. All four indices coincide, so $v=1$ and the contribution is $N^{1-3}=N^{-2}$, against $N^{-1}$ in the real case. The two non-crossing pair partitions are unaffected and still contribute $1$ each.

\Q{2.18}{A} The histogram matches $\rho_{sc}$ visibly by $N=2000$, and the $\pm1/\sqrt N$ picture is the same: the limit depends on the entries only through their variance, although Wick is false for them.

\end{document}
